%% Experimental results moved out of \Cref{sec:exp} to keep the main text short.  Nothing here is
%% duplicated in the main text.  Shared by ../paper/main.tex and ../submission/body.tex.

\subsection{The two designs whose point estimate has not converged}
\label{app:hard}

The Markov and Logistic columns need a caveat of their own, and it is the sharpest test of
whether our variance estimator is really doing the work we claim. Neither design's point estimate
has reached its asymptotic regime at $N=200{,}000$, for two different reasons. The Markov design's
regime chain has an integrated autocorrelation time of $\HardTau$ observations --- the sum
$1+2\sum_{k\ge1}$ of its autocorrelations, which is the factor by which dependence multiplies the
variance of a sample mean, so that $N$ dependent observations carry about as much information as
$N/\HardTau$ independent ones --- so its effective sample size is a small fraction of $N$ and the
point estimate is $\HardRmseSmall$ times the efficient standard deviation. On the Logistic design
the cause is the curvature rather than the dependence: the weight $\alpha=\sigma'(x^{\!\top}\theta)$
is largest at $\theta_0=0$, where the warm-up evaluates it, so $\Hb$ built there overestimates $H$
and the early steps are correspondingly too short; the point estimate is $\HardLogitRmseSmall$
times efficient and the fitted standard error is biased low with it. An interval built from the
\emph{correct} asymptotic variance must undercover in that situation, no matter who supplied
the variance, simply because $\sqrt N(\theta_T-\ts)$ is not yet close to its limit. So on these two designs the
question is not whether coverage is nominal but whether our interval matches the interval that
uses the exact $H^{-1}SH^{-1}$ at the same iterate. \Cref{tab:hard} and \Cref{fig:hard} answer it:
over a sweep in $N$ both rise towards nominal together while the dependence-blind interval does
not. On the Markov design the agreement is essentially exact: our coverage is within
$\HardMaxGapMarkov$ of the oracle's at every $N$, and at the top of the sweep, $N=\HardNBig$, both are
$\HardCovOursBig$ and $\HardCovOracleBig$ with the point estimate down to $\HardRmseBig$ times
efficient, against $\HardCovPluginBig$ for the dependence-blind interval. The direction of travel
is the whole point. Across the fivefold increase in $N$ our coverage and the oracle's move
together towards $0.95$ and stay within $\HardMaxGapMarkov$ of each other, while the
dependence-blind interval creeps up by a few hundredths towards its own asymptote, which
\Cref{prop:coverage} puts at $2\Phi(z_{0.025}/\sqrt{\KappaLoMarkov})-1$, far below nominal. We
stop the sweep at $N=\HardNBig$ rather than continuing to where coverage is nominal, because what
the experiment is for is the gap between the first two columns, and \Cref{tab:hard} shows that gap
is already at the level of the Monte Carlo standard error there. On the Logistic design
there is a real gap, and it is ours: our coverage is $\HardLogitCovOursSmall$ against the
oracle's $\HardLogitCovOracleSmall$ at $N=\HardNSmall$, a difference of $\HardLogitGapSmall$,
falling monotonically to $\HardLogitGapBig$ at $N=\HardNBig$, the top of the sweep. The cause is
visible in \Cref{tab:main}: our fitted width there is $\WidthRatioLogisticAbs$ times the oracle's,
i.e.\ the standard error is biased low, because $\Hb$ is an average over the whole stream and the
early blocks contribute $\alpha=\sigma'(0)=1/4$, its largest possible value, so $\Hb$
overestimates $H$ and $\Hb^{-1}S\Hb^{-1}$ underestimates the sandwich. That bias decays with the
Cesàro average of $\norm{\theta_t-\ts}$, which is why the gap closes; it is a finite-sample
property of averaging curvature along the path, it affects the base method identically, and no
choice of long-run covariance estimator would remove it. That is the signature we want: our error tracks
the oracle's at every sample size, and the dependence-blind error is the one that survives
$N\to\infty$. The step-length safeguard bound at most $\HardClipMax$ times per run anywhere in
that sweep, and the count does not rise with $N$, which is the empirical content of
\Cref{prop:shift}.

We report every design in \Cref{tab:main} at the same $N$ rather than tuning $N$ per design,
which is why the Markov column looks worst there; \Cref{tab:hard} is where that column is
resolved.

\paragraph{The step-length safeguard is what makes the hard design run at all.}
\Cref{tab:ablation}(i) varies the safeguard. With none, the autoregressive designs are barely
affected ($\ShiftEffArhomNone$ against $\ShiftEffArhomCap$ in relative variance) but the
regime-switching design diverges: relative variance $\ShiftEffMarkovNone$ against
$\ShiftEffMarkovCap$ with the safeguard, and coverage $\ShiftCovMarkovNone$ against
$\ShiftCovMarkovCap$. The safeguard is not doing this by damping the algorithm into submission:
it binds a median of $\ShiftClipMarkovCap$ times in $10{,}000$ blocks on that design and
$\ShiftClipArhomCap$ times on AR-hom, and \Cref{tab:hard} shows the count flat as $N$ grows by a
factor of five.

The constant $c_D$ behaves the way a safeguard constant should, which is to say it is irrelevant
until it is not. On AR-hom, changing it by a factor of four in either direction moves nothing
($\ShiftEffArhomCapLo$ at $c_D=1/4$ and $\ShiftEffArhomCapHi$ at $c_D=4$, against
$\ShiftEffArhomCap$ at our default $c_D=1$). On the design that actually needs the safeguard both
directions cost something, and asymmetrically: $c_D=1/4$ gives $\ShiftEffMarkovCapLo$ against
$\ShiftEffMarkovCap$ --- binding $\ShiftClipMarkovCapLo$ times instead of
$\ShiftClipMarkovCap$, it now restrains steps that were fine --- while $c_D=4$ is too loose to
prevent the blow-up at all and the relative variance is $\ShiftEffMarkovCapHi$. So $c_D$ is a
tuning constant with a real failure mode on one side, and we say so rather than advertising
insensitivity we only observe on the easy designs.

The schedule shift of \Cref{app:t0} is the alternative the contraction condition suggests more
directly, and it is measurably worse: $\ShiftEffMarkovShift$ against $\ShiftEffMarkovCap$ on the
regime-switching design, because it slows all $10^4$ steps to correct a handful. We report it
because we implemented it first, and because the comparison is the evidence for the design we
kept.

\subsection{The first-order baselines and the conditioning of the design}
\label{app:firstorder}

The first-order rows of \Cref{tab:main} need a caveat, and it cuts against us if we hide it.
These designs are ill-conditioned --- $\mathrm{cond}(H)=\CondArhom$ for AR-hom --- which is precisely
the regime streaming second-order methods exist for, and there averaged SGD's \emph{point}
estimate has not converged: its root-mean-square error is $\RmseAsgdArhom$ times the
efficient value against $\RmseBgsnArhom$ for \bgsn{}. Its intervals are therefore
uninformative about the inference question, whatever covariance they use. Two pieces of evidence show that this is about conditioning and not about our construction.
First, within \Cref{tab:main} itself: streaming AdaGrad, which rescales per coordinate, is
efficient on AR-het ($\RmseAdagradArhet$) and with our long-run variance it covers
$\CovAdagradArhet$ --- so the construction transfers to a first-order method as soon as that
method's point estimate converges. Second, \Cref{tab:wellcond} repeats the AR-hom design with
$\Sigma_x=I$, where every point estimator is efficient ($\WcRmseAsgd$ for averaged SGD,
$\WcRmseBgsn$ for \bgsn{}) and only the covariance can decide coverage: averaged SGD with our
long-run variance covers $\WcCovAsgdOurs$, the same estimator with the i.i.d.\ plug-in covers
$\WcCovAsgdPlugin$, and \bgsn{} covers $\WcCovBgsn$.

\subsection{Distributional accuracy beyond a single quantile}
\label{app:ks}

Coverage checks a single quantile, while the theorem claims a whole distribution, so we also
report the Kolmogorov--Smirnov distance between the studentised error (the error divided by
its own estimated standard error, which is what an interval implicitly compares to a normal
quantile)
$\sqrt N(\theta_{T,j}-\ts_j)/\{u_j^{\!\top}\Sft u_j\}^{1/2}$ and the standard normal, taken
per coordinate across replicates. On AR-hom the worst of the $d$ coordinates gives $\KsBgsnArhom$ against a $5\%$ critical value of
$\KsCritArhom$: marginally outside, so at this number of replicates the studentised distribution
is distinguishable from normal, and we report that rather than rounding it in our favour. The
dependence-blind studentisation gives $\KsPluginArhom$, an order of magnitude further out. The
comparison, not the test outcome, is the content: our studentisation is at the edge of detectable
non-normality while the dependence-blind one is not close. The standard errors themselves are accurate: $\SeRatioBgsnArhom$ times the
truth on average.

\subsection{Our estimator against an iterate-path estimator}
\label{app:obm}

Panel (a) puts our estimator and an iterate-path estimator on the same object, because an
interval uses the whole sandwich $\Sigma=H^{-1}SH^{-1}$ rather than $S$ alone. We compare our
composite $\Hb^{-1}\Sft\Hb^{-1}$, whose relative error falls with fitted exponent
$\RateExpSandwich$, with overlapping batch means on the \emph{iterate} path at the block length
its analysis prescribes --- the estimator of $\Sigma$ that uses no gradients and no curvature.

Two things must be said carefully here, because the comparison is easy to overstate. First, the
iterate-path estimator's relative error is flat over our range: the fitted exponent is
$\RateExpIterateObm$, and we do \emph{not} read that as its rate. Its published rate is
$N^{-1/8}$, which over a sixteenfold increase in $N$ predicts a reduction by a factor
$16^{-1/8}=0.71$ --- a change small enough that our range cannot separate it from zero, and our
fitted exponent should be treated as uninformative about it rather than as a measurement.
Second, what our range \emph{does} establish is the level, and the level is the practical
statement: over the whole range the iterate-path estimator's relative error in $\Sigma$ sits
near $\RateObmLevel$, which is not a usable standard error whatever its asymptotic rate, while
ours falls from $\RateOursLevelSmall$ to $\RateOursLevelBig$. We are comparing an estimator that
has entered its asymptotic regime with one that has not, and the honest claim is about the
sample sizes a practitioner would actually have, not about which exponent is larger.

\subsection{Curvature schedules}
\label{app:sched}

\RefGapFig{} compares deterministic gapping with Bernoulli thinning at a matched number of
rank-one updates, with the optimisation frozen at $\ts$ so that only estimation quality is
measured. The excess variance over the independent-data benchmark follows the two predicted
curves --- $O(\rho^{m})$ for gapping, $(\kappa_H-1)/m$ for thinning --- over more than an
order of magnitude in $m$. As stated in \Cref{sec:limits}, this is a curvature-side result:
\Cref{tab:ablation}(f) shows that the end-to-end effect on coverage is within Monte Carlo
error, which is what the theory predicts, since $\Hb$ enters the limit only through a rate.

Where gapping does pay is cost, and \Cref{tab:main} measures it. Compare \bgsn{} with the row
``base method with our long-run variance'': the two use the same variance estimator and the same
number of rank-one curvature updates per observation ($1/d$), and they cover the same
($\CovBgsnArhom$ against $\CovSniidtsArhom$ on AR-hom). They differ only in \emph{how} those
updates are placed --- on a deterministic gap of $m=d$ against Bernoulli retention with $p=1/d$.
At these settings the wall-clock difference is $\GapSpeedupArhom\times$, which is within our
timing noise, so we do not claim a speed-up: a gap avoids drawing and testing a random variable
per observation, but that saving is small next to the $O(d^2)$ update it guards. The case for
gapping is therefore the one \Cref{prop:gap} makes and \RefGapFig{} measures --- a strictly
better curvature estimate at a matched number of updates, $O(\rho^m)$ excess variance against
$(\kappa_H-1)/m$ --- together with the fact that it makes the per-block cost deterministic
rather than random. It is a better cost knob, not a faster one.
\subsection{Ablations}
\label{app:abl}

\RefAblFig{} and \Cref{tab:ablation} vary one design choice at a time.
(a) Coverage rises monotonically to nominal with $N$, and the gap to the dependence-blind
interval does \emph{not} close --- it converges to the predicted $0.749$.
(b) Coverage as a function of $b$ is flat for the two-scale estimator over
$b\in[10,320]$ and degrades at both ends for plain batch means (bias at small $b$, noise at
large $b$): the practical content of the $O(\rho^{b})$-versus-$O(1/b)$ distinction.
(c) The curvature warm-up matters, though less than it did before the step-length safeguard
absorbed part of its job: with essentially no warm-up, coverage on this design is
$\CovNoWarmup$ against $\CovBgsnArhom$ at the default, and above $c_{\mathrm w}\approx25$ the
result is insensitive. (c$'$) is where the warm-up is load-bearing: on the regime-switching
design the relative variance is $\WarmMarkovFifty$ at $c_{\mathrm w}=50$ against
$\WarmMarkovTwoHundred$ at $c_{\mathrm w}=200$, because there $50d$ consecutive observations span
only about thirty regime visits and leave the curvature estimate rank-starved. That is what sets
our default, and it is a statement about effectively distinct design points rather than about
observations.
(d) Results are insensitive to the ridge exponent $\qr$ over the whole admissible range, and to
the ridge constant $c_\iota$ from $0$ --- no ridge at all --- up to our default $0.01$; at
$c_\iota=0.1$ coverage slips by two points and at $c_\iota=1$ the ridge dominates the data terms
and the estimator is no longer efficient. So the ridge is safe to leave alone but not safe to
enlarge, and the honest reading is that the default sits at the top of the flat region rather than
in the middle of it. That $c_\iota=0$ works as well as the default is worth stating plainly: the
ridge is in the algorithm so that \Cref{lem:ridge} can hold for \emph{every} $t$ without an
assumption on the iterates, not because the runs need it.
(e) Results are insensitive to the initial curvature scale $\lambda_0$ over eight orders of
magnitude, which matters because that is a parameter a user would otherwise have to think about.
(f) Gap and thinning behave as \Cref{prop:gap} predicts and do not move coverage, which is
what the theory says: $\Hb$ enters the limit only through a rate.
(g) At strong dependence ($\phi=\psi=0.85$, $\kappa=\KappaStrong$) the two estimators separate in
exactly the way the $O(\rho^b)$-versus-$O(1/b)$ distinction predicts. Plain batch means need a
long block: coverage rises from $\StrongCovBmSmallB$ at $b=20$ to $\StrongCovBmBigB$ at $b=160$.
The two-scale correction does not: it is already at $\StrongCovFtSmallB$ at $b=20$ and stays there
($\StrongCovFtBigB$ at $b=160$). The free diagnostic $r_j$ falls from $\StrongAdeqSmallB$ to
$\StrongAdeqBigB$ over the same range, which is what makes it usable as a guide to whether plain
batch means would have been adequate. So the correction is what buys the short block, and the short
block is what keeps the estimator cheap.
(h) The weighted averaged variant is slightly better at small $N$ and slightly worse at
large $N$; both are reported.
The non-positive-definiteness fallback of \Cref{rem:psd} never triggered anywhere: the rate is
$0$ in every simulation and $\RealPsdFallback$ of coordinate-replicate pairs on the real streams,
where an earlier version of these experiments did not record it at all and where we would have
expected it if anywhere, since there the lag-one correction is largest. It never triggered
($0.0\%$ of all runs, all panels).
\subsection{Cost: the correct interval is also the cheaper one}
\label{app:costexp}

The authoritative statement is the operation count, because it is exact and independent of the
machine. The streaming pass grows as $d^{\StreamExpBgsn}$ for \bgsn{} and as
$d^{\StreamExpPlugin}$ once the i.i.d.\ plug-in variance is added, because the plug-in score
covariance costs $O(d^2)$ per \emph{observation} whereas our long-run covariance costs $O(d^2)$
per \emph{block}; evaluating the whole count at $N=10^{8}$ gives $d^{\AsymExpBgsn}$ against
$d^{\AsymExpPlugin}$. Those are the design claim, and they are arithmetic rather than a fit.

The exponents fitted at the $N$ of \Cref{tab:cost} are \emph{not} the streaming cost, and we
say so rather than quoting them as if they were: with $c_{\mathrm w}=200$ and $N=200{,}000$ the
one-off $O(c_{\mathrm w}d^3)$ warm-up dominates the $O(dN)$ pass for $d\ge40$
(the crossover is $N\gtrsim c_{\mathrm w}d^2$), so both measured exponents sit near $2$. What
is robust at this $N$ is the \emph{level}: the plug-in variant costs more than \bgsn{} in
wall-clock at every $d$ we tried, by a factor between $\CostRatioPluginLo$ and
$\CostRatioPluginHi$ (median $\CostRatioPlugin$), and in the main configuration \bgsn{}
runs in $\TimeRatioArhom$ times the time of the dependence-blind streaming Newton method --- that
is, our valid intervals cost less than the invalid ones, not more. The measured wall-clock
exponents, on a shared and heavily loaded machine, are $d^{\CostExpBgsn}$ and
$d^{\CostExpPlugin}$; they are noisier than the counts and we report them only as corroboration.
\RefCostFig{} plots the same timings against $d$, where the message is the vertical separation
rather than either slope.

\subsection{Choosing the block length on a real stream}
\label{app:realb}

\Cref{fig:real}(c) shows why the block length is the one parameter a practitioner must think
about on real data, and that the free diagnostic
$r_j=|u_j^{\!\top}(\Sft-\Sbm)u_j|/(u_j^{\!\top}\Sbm u_j)$ flags the problem. The denominator is
the batch-means quadratic form, not the two-scale one: $\Sbm$ is a sum of outer products and so
its quadratic forms are non-negative, whereas those of $\Sft$ can be negative
(\Cref{rem:psd}), and an earlier version of this diagnostic divided by the latter and produced
meaningless numbers on the real streams. On the traffic stream $r_j$ is
$\RealAdeqMetro$ at the $b=\RealBMetro$ we use, and the sweep in \Cref{fig:real}(c) shows what it is
good for: $r_j$ is U-shaped in $b$, rising both for short blocks, where plain batch means carry the
$O(1/b)$ bias, and for long ones, where too few blocks remain for a stable estimate. Our coverage
over that fortyfold range in $b$ is flat and stays within a few hundredths of the infeasible
oracle-variance curve, which is the evidence that the shortfall from nominal on these segments is
the point estimate rather than the covariance. So $r_j$ is a guide to choosing $b$, not a
calibration guarantee, and we do not claim it predicts coverage.
